\section*{THE FORMALISATION OF LOGIC.}
\phantomsection
\addcontentsline{toc}{section}{THE FORMALISATION OF LOGIC.}

The general impression which seems to emerge from the (very patchy) texts that we have on Greek philosophical thought of the Vth century, is that it is dominated by an increasingly conscious effort to extend to the whole field of human thought the procedures for conducting discussion put in hand with so much success by contemporary rhetoric and mathematics - in other words, to create Logic in its most general form. The tone of philosophical writings undergoes at this time a sudden change: while in the VIIth and VIth centuries the philosophers affirm or vaticinate (or at most outline vague arguments, based on equally vague analogies), starting with Parmenides and especially Zenon, they argue and try to draw out general principles which can serve as a basis for their dialectic: it is in Parmenides that one finds the first statement of the principle of the excluded middle, and the proofs ``by contradiction'' of Zenon of Elea remain famous. But Zenon writes in the middle of the Vth century; and, whatever the uncertainties of our documentation, \(^{2}\) it is very likely that at this time, mathematicians, in their own sphere, were currently using these principles.

As we said above, it is not for us to retrace the innumerable difficulties which abound at every step in the gestation of this Logic, and the controversies which result, from the Eleatics to Plato and Aristotle, via the Sophists; let us pick up here only the role played in this evolution by the assiduous cultivation of oratorical art and the analysis of language which is one of its corollaries, developments that it is agreed should be attributed mainly to the Sophists of the Vth century. On the other hand, even if the influence of mathematics is not always explicitly recognised, it is none the less manifest, in particular in the writings of Plato and Aristotle. It can be said that Plato was almost obsessed by mathematics; without being himself an inventor in this area, he kept up, after a particular period in his life, with the discoveries by contemporary mathematicians (many of whom were his friends or his pupils), and never after ceased from a most direct kind of interest, even going so far as to suggest new directions of research; so it is that constantly, in his writings, mathematics serves as illustration or model (and sometimes even, as with the Pythagoreans, feeding his leanings towards mysticism). As for his pupil Aristotle, he could not have avoided receiving the minimum of mathematical foundations required of the pupils of the Academy, and a volume has been produced consisting of extracts of his work that are about mathematics or refer to it {[}153 d{]}; but he seems never to have made a great effort to keep in touch with the mathematical movements of his time, and he only quotes in this area results that had been popularised a long time before. This displacement will besides only become more marked with the majority of later philosophers, of whom many, due to the lack of technical preparation, believe themselves in all good faith to be speaking knowledgeably about mathematics, whereas they will only be referring to a stage long since overtaken in its evolution.

The end product of this period, as far as Logic is concerned, is the monumental work of Aristotle {[}6{]}, whose great merit is to have succeeded in systematising and codifying for the first time procedures of reasoning which had remained vague or unformulated in his predecessors. \(^{3}\) We must above all retain here, as our aim, the general thesis of this work, to know that it is possible to reduce all correct reasoning to the systematic application of a small number of immutable rules, which are independent of the particular nature of the objects in question (an independence which is clearly demonstrated by the representation of the concepts or propositions by means of letters --- probably borrowed by Aristotle from the mathematicians). But Aristotle concentrates his attention almost exclusively on a particular type of relations and logical chains, making up what he calls the ``syllogism'': it is mainly a case of relations that we would translate nowadays in the form \(A \subset B\) or \(A \cap B \neq \emptyset\) in the language of the theory of sets, \(^{4}\) and with the way of chaining these relations or their negations, by means of the schema

\[
(A \subset B \text {   and   } B \subset C) \Rightarrow (A \subset C).
\]

Aristotle was still sufficiently informed of the mathematics of his era in order to have seen that schemas of this type were not sufficient to take account of all the logical operations used by mathematicians, nor, all the more, of the other applications of Logic ({[}6{]}, An. Pr., I, 46). \(^{5}\) At least the deep study of the different forms of ``syllogisms'' which he undertakes (and which is almost entirely consecrated to the elucidation of the perpetual difficulties raised by the ambiguity or the obscurity of the objects on which the reasoning bears) gives him among others the chance to formulate the rules for obtaining the negation of a proposition ({[}6{]}, An. Pr., I, 46). It is also to Aristotle that is due the credit for having distinguished with great precision the role of ``universal'' propositions from that of ``particular'' propositions, the first sketch of quantifiers. \(^{6}\) But it is too well known how the influence of his writings (often interpreted in a narrow and unintelligent way), which remain still very influential until well into the XIXth century, was to encourage philosophers in their neglect of the study of mathematics, and to block the progress of formal Logic. \(^{7}\)

However this latter continued to make progress in Antiquity, in the surroundings of the Megaric and Stoic schools, rivals of the Peripatetics. Our information about these tenets are unfortunately all at second hand, often passed on by adversaries or mediocre commentators. The essential progress achieved by these logicians consists, it would appear, of the formation of a ``propositional calculus'' in the meaning which it has today: instead of being restricted, like Aristotle, to propositions of the particular form \(A \subset B\) , they state rules about completely indeterminate propositions. Further, they had analysed the logical links between these rules in such a deep way that they knew how to deduce them all from five of them, set down as ``unprovable'', by means of procedures very similar to modern methods {[}23{]}. Unfortunately their influence was fairly ephemeral, and their results were to fall into oblivion until the day when they were rediscovered by the logicians of the XIXth century. Aristotle remains as the uncontested master in Logic until the XVIIth century; it is known in particular that the scholastic philosophers are entirely under his sway, and if their contribution to formal logic is far from negligible {[}25{]}, it does not contain any progress at the highest level compared with the achievements of the philosophers of Antiquity.

It is appropriate, however, to note here that it does not seem that the works of Aristotle or his successors had a notable influence on mathematics. The Greek mathematicians conducted their research along the path opened by the Pythagoreans and their successors of the IVth century (Theodorus, Theeta, Eudoxus) without apparently bothering with formal logic in the presentation of their results: a finding that should hardly astonish when one compares the flexibility and precision acquired already by that time by mathematical reasoning, to the very rudimentary state of Aristotelian logic. And when logic will overtake this stage, it is again the new acquisitions from mathematics that will guide it in its evolution.

With the development of algebra, it was indeed impossible to avoid being struck by the analogy between the rules of formal logic and the rules of algebra, the ones like the others having in common the property of being applicable to objects (propositions or numbers) which are not precisely determined. And when in the XVIIth century algebraic notation took its definitive form in the hands of Viète and Descartes, almost immediately one can see appearing diverse attempts at a symbolic notation intended for the representation of logical operations; but before Leibniz, these tentative efforts, for example like that of Herigone (1644) at writing down the proofs of elementary Geometry, or that of Pell (1659) at writing down those of Arithmetic, remain very superficial and do not lead to any progress in the analysis of mathematical reasoning.

With Leibniz, we are in the presence of a philosopher who is also a mathematician of the first rank, and who will know how to draw out of his mathematical experience the germ of the ideas that will bring formal logic out of the scholastic dead-end. \(^{8}\) A universal spirit if ever there were one, inexhaustible source of original and fruitful ideas, Leibniz would be interested even more in Logic as it found itself at the very heart of his grand projects for formalising language and thought, at which he never ceased working throughout his life. Trained in his childhood in scholastic logic, he was seduced by the idea (going back to Raymond Lulle) of a method which would reduce all human concepts to primitive concepts, making up an ``Alphabet of human thought'', and would rebuild them in a quasi mechanical way to obtain all true propositions ({[}198 b{]}, v. VII, p.~185; cf.~{[}69a{]}, chap.~II). Still very young, he had also conceived another much more original idea, that of the usefulness of symbolic notations as an ``Ariadne's string'' string of thought: \(^{9}\) ``The true method'', he said, ``must provide us with a filum Ariadnes, that is to say a kind of sensitive and coarse means that guides the mind, in the same way as lines drawn in geometry and the type of operations that are prescribed to apprentices in Arithmetic. Without that our mind would not know how to go along a long path without straying.'' ({[}198 b{]}, v. VII, p.~22; cf.~{[}69 a{]}, p.~90). Knowing little of the mathematics of his time until about his 25th year, it is at first in the form of a ``universal language'' that he presents his projects ({[}69 a{]}, chap.~III); but as soon as he comes in contact with Algebra, he adopts it as model for his ``universal Characteristic''. By that he means a type of symbolic language, capable of expressing without ambiguity all human thought, of reinforcing our powers of reasoning, of avoiding errors by means of an effort of entirely mechanical concentration, finally constructed so that ``the chimera that the person himself who presents them does not hear, would be impossible to write down in these characters'' ({[}198 a{]}, v. I, p.~187). In innumerable passages in his writings where Leibniz alludes to this grandiose project and to the progress to which its attainment would lead (cf.~{[}69 a{]}, chap.~IV and VI), it is seen with what clarity he had conceived the notion of a formalised language, a pure combination of symbols of which only the chaining is important, \(^{10}\) in such a way that a machine would be able to produce all the theorems, \(^{11}\) and that all controversies would be resolved by a simple calculation ({[}198 b{]}, v. VII, p.~198-203). If these hopes might appear excessive, it is none the less true that it is to this constant tendency of Leibniz's thought that must be tied a good part of his mathematical work, beginning with his work on the symbolism of the infinitesimal Calculus (see p.~190 ff.); he was himself perfectly aware of this, and linked explicitly also to his ``Characteristic'' his ideas on indicial notation and determinants ({[}198 a{]}, v. II, p.~204; cf.~{[}69 a{]}, pp.~481-487) and his sketch of ``geometrical Calculus'' (see pp.~50 and 61 ff.; cf.~{[}69 a{]}, chap IX). But in his mind the essential part had to be symbolic Logic, or, as he puts it a ``Calculus ratiocinator'', and though he does not succeed in creating this calculus, at least we see him trying to at least three times. On his first attempt, he has the idea of associating with each ``primitive'' term a prime number, each term made up of several primitive terms being represented by the product of the corresponding prime numbers; \(^{12}\) he seeks to translate into this system the usual rules of syllogism, but runs up against major complications caused by negation (that he tries, fairly naturally, to represent by a change of sign) and abandons rapidly this path ({[}198 c{]}, pp.~42-96; cf.~{[}69 a{]}, pp.326-344). In later attempts, he seeks to give Aristotelian logic a more algebraic form; sometimes he keeps the notation \(AB\) for the conjunction of two concepts; sometimes he uses the notation \(A + B\) ; \(^{13}\) he notes (in multiplicative notation) the law of idempotence \(AA = A\) , remarks that one can substitute the proposition ``all A is B'' by the equality \(A = AB\) and that one can recover starting from there most of the rules of Aristotle by a purely algebraic calculation ({[}198 c{]}, pp.~229-237 and 356-399; cf.~{[}69 a{]}, pp.345-364); he also has the idea of the empty set (``non Ens''), and recognises for instance the equivalence of the propositions ``all A is B'' and ``A.(not B) is not'' (loc. cit.). Further, he remarks that his logical calculus is applicable not only to the logic of concepts, but also to that of propositions ({[}198 c{]}, p.~377). He appears to be therefore very near to ``Boolean calculus''. Unfortunately, it appears that he was not able to free himself completely from the scholastic influence; not only does he set himself as almost the only aim of his calculus the transcription, into his notation, of the rules of syllogism, \(^{14}\) but he goes as far as sacrificing his most felicitous ideas to the desire to recover completely the rules of Aristotle, even those that were incompatible with the idea of an empty set. \(^{15}\)

The work of Leibniz remained for the most part unpublished until the beginning of the XXth century, and had only little direct influence. During the whole of the XVIIIth and the beginning of the XIXth centuries, different authors (de Segner, J. Lambert, Ploucquet, Holland, De Castillon, Gregonne) sketched attempts similar to those of Leibniz, without ever progressing substantially beyond the point at which he had stopped; their work only had a feeble effect, which means that most of them ignored all results due to their predecessors. \(^{16}\) It is in any case under the same conditions that G. Boole, who must be considered to be the real creator of modern symbolic logic {[}29{]}, writes. His key idea consists in putting himself systematically in the position of considering ``extension'', so of calculating directly with sets, and writing xy for the intersection of two sets, and \(x + y\) for their union when x and y have no element in common. He introduces as well a ``universe'' denoted 1 (the set of all elements) and the empty set denoted 0, and he writes 1 - x for the complement of x. As Leibniz had done, he interprets the relation of inclusion by the relation xy = x (from which he extracts without difficulty the justification of the rules of classical syllogism) and his notation for union and complements gives his system a flexibility that was missing in his forerunners. \(^{17}\) Furthermore, by associating with each proposition the set of ``cases'' in which it holds, he interprets the relation of implication as an inclusion, and his calculus of sets gives him in this way the rules of the ``propositional calculus''.

In the second half of the XIXth century, Boole's system formed the basis for the work of an active school of logicians who improved it and completed it at several points. It is thus that Jevons (1864) enlarged the meaning of the operation of union \(x + y\) by extending it to the case in which \(x\) and \(y\) are arbitrary; A. de Morgan in 1858 and C. S. Peirce in 1867 proved the duality laws

\[
(\mathbf {C} A) \cap (\mathbf {C} B) = \mathbf {C} (A \cup B), (\mathbf {C} A) \cup (\mathbf {C} B) = \mathbf {C} (A \cap B); ^{18}
\]

De Morgan tackles also, in 1860, the study of relations, defining inversion and the composition of binary relations (that is to say the operations corresponding to the operations \(\overline{G}^{-1}\) and \(G_{1} \circ G_{2}\) on graphs). \(^{19}\) All this work is systematically expounded and developed in the massive and prolific work of Schröder {[}277{]}. But it is fairly interesting to note that the logicians of whom we have been speaking do not seem very interested in the application of their results to mathematics and that, on the contrary, Boole and Schröder especially seem to have as their principal aim to develop ``Boolean'' algebra by imitating the methods and problems of classical algebra (often in a fairly artificial way). The reasons for this attitude must doubtless be seen in the fact that the Boolean calculus still lacked a facility for transcribing most of mathematical reasoning, \(^{20}\) and only supplied in this way a very partial answer to the great dream of Leibniz. The construction of a formalism better adapted to mathematics --- of which the introduction of variables and quantifiers, due independently to Frege {[}117 a, b, c{]} and C. S. Peirce {[}248 b{]}, make up the major stage --- was the work of logicians and mathematicians who, unlike the above, had above all as their aim applications to the foundations of mathematics.

Frege's project {[}117 b and c{]} was to create a foundation for arithmetic based on a logic formalised by a ``writing of concepts'' (Begriffschrift) and we will come back later (p.~29) to the way in which he defines the natural numbers. His work is characterised by extreme precision and attention to detail in the analysis of concepts; it is because of this tendency that he introduces many a distinction which turns out to be of great importance in modern logic: for instance it is he who first distinguishes between the statement of a proposition and the assertion that this proposition is true, between the relation of belonging and that of inclusion, between an object x and the set \(\{x\}\) reduced to this single object, etc.. His formalised logic, which contains not only ``variables'' with the meaning used in mathematics, but also ``propositional variables'' representing indeterminate relations, susceptible to quantification, would later (through the work of Russell and Whitehead) supply the fundamental tool for metamathematics. Unfortunately, the symbols which he adopts are hardly inspiring, of a terrifying typographic complexity and far removed from the practice of mathematicians; which had the result of turning away these latter and reducing considerably the influence of Frege on his contemporaries.

Peano's goal was at the same time much vaster and much more down to earth; it consisted in publishing a ``Formulary of mathematics'', written entirely in formal language and containing, not only mathematical logic, but all the results of the most important branches of mathematics. The speed with which he managed to complete this ambitious project, helped by a host of enthusiastic collaborators (Vailati, Pieri, Padoa, Vacca, Vivanti, Fano, Burali-Forti) is witness to the excellence of the symbolism which he had adopted: following closely the current practice of mathematicians, and introducing numerous well-chosen abbreviating symbols, his language remains as well fairly easily legible, thanks notably to an ingenious system for replacing brackets by full stop separators {[}246 f{]}. Much notation due to Peano is today adopted by the majority of mathematicians: we quote \(\in, \supset\) (but, contrary to the present use, with the meaning of ``is contained in'' or ``implies'' \(^{21}\) ), \(\cup, \cap, A - B\) (set of differences \(a - b\) where \(a \in A\) and \(b \in B\) ). On the other hand, it is in the ``Formulary'' that is found for the first time a thorough analysis of the general notion of function, of those of direct image \(^{22}\) and reciprocal image, and the remark that a sequence is only a function defined on N. But quantification, with Peano, is subject to hampering restrictions (one can only, in principle, quantify, in his system, relations of the form \(A \to B\) , \(A \Leftrightarrow B\) or \(A = B\) ). Further the almost fanatical zeal of some of his disciples laid them wide open to ridicule; criticism, often unjustified, by H. Poincaré in particular, was a heavy blow to the Peano school and became an obstacle to the diffusion of his doctrines in the mathematical world.

With Frege and Peano the essential elements of the formal languages used today were acquired. The most widespread is doubtless that hammered out by Russell and Whitehead in their great work ``Principia Mathematica'', which happily links the precision of Frege and the convenience of Peano {[}266{]}. Most actual formal languages are differentiated from it only by changes of secondary importance, aimed at simplifying its use. Among the most ingenious, we quote the ``functional'' writing of relations (for instance \(\in xy\) instead of \(x \in y\) , thought up by Lukasiewicz, thanks to which brackets can be completely omitted; but the most interesting is without doubt the introduction by Hilbert of the symbol \(\tau\) , which allows the consideration of the quantifiers \(\exists\) and \(\forall\) as abbreviation signs, the avoidance of the introduction of the ``universal'' functional symbol \(\iota\) of Peano and Russell (which is only applied to functional relations), and finally which avoids the need to formulate the axiom of choice in the theory of sets ({[}163 a{]}, v.III, p.~183).

